% Quantitative source calculation for the legacy weighted response limit.
\section{A priced mission forces the weighted response to vanish}
\label{subsec:ex-trace-small-means}

This section records the quantitative argument behind the weighted
field-image estimate used by the trace receiver. It exposes the local price,
scale, and detector inputs. It does not derive the mission's local price
bound from the finite-volume estimates.

Let $S$ be the selected interacting table, and let
$\mathcal U_i^{(d)}:E_d\to S_d$ be its scale chart from the reference grade
space $E_d$ into the selected grade space $S_d$. Here $E_d$ is the
grade-$d$ field space of the reference (standard) table; it carries a
complex-antilinear real-structure involution and an action $R(s)$ of the
signed permutations $s$ of the four axes. Let $\mathcal V$ be the
finite-dimensional real vector space of grade-five response arrays
$y=(y_\alpha)_{\alpha\in\mathrm{Axis}}$, $y_\alpha\in E_5$, which are
real and signed-equivariant: every $y_\alpha$ is fixed by the real-structure
involution, and
\[
 R(s)y_\alpha=\operatorname{sign}(s,\alpha)\,y_{s\alpha}
 \qquad\text{for every signed permutation }s\text{ and every axis }\alpha,
\]
where $\operatorname{sign}(s,\alpha)$ is the sign that $s$ attaches to
the axis $\alpha$ (this is the source's one-point response table). Both
conditions are part of the definition; signed equivariance alone defines a
larger space. Let $\mathcal B$ be the real vector space of arrays of raw
grade-five labels defined by the same two conditions, with the conjugation
and the signed-permutation action of the raw labels in place of those of
$E_5$. The normalized response map is a real-linear map
$s_i:\mathcal B\to\mathcal V$. Write $j_i$ for the
selected frame level and $k_i$ for its associated scale index, the summit
index of $j_i$; in particular $k_i>j_i$, and $j_i\to\infty$ along the
selected sequence. The source
uses two distinct positive scales:
\[
 z_i=\widehat\beta_{k_i},\qquad \zeta_i=Z.\beta(k_i).
\]
The first, $z_i$, appears in the polynomial weight and rate estimate; the
second, $\zeta_i$, is the Arch scale in the medium and cancels from the
response identity below. Let $a_i$ be the field dilation at level $j_i$,
$b_i$ the frame scale, $m_i$ the frame cap, and $r_i$ the physical dilation.
The selected scale relations are
\[
 r_i=a_ib_i,\qquad 0<b_i\leq m_i,
 \qquad m_i=\operatorname{tiltcap}(j_i).
\]
Let $v_i\in\mathcal B$ be the selected medium and $x_i\in\mathcal B$ its
core field. The three identities
\eqref{eq:ex-trace-medium-core}--\eqref{eq:ex-trace-night-scaling} below
are used only for all sufficiently large $i$, and the source supplies the
first two only on such a tail; the third requires $k_i$ to be at least the
grade-five reference level, which holds for large $i$ because
$k_i>j_i\to\infty$. The core
representation is
\begin{equation}
 v_i(\alpha)=\zeta_i^{-1}\mathcal E_{k_i}(x_i(\alpha))
 \qquad(\alpha\in\mathrm{Axis}),
 \label{eq:ex-trace-medium-core}
\end{equation}
where $\mathcal E_k$ is the source's core-to-strip map. The response
normalization is, for every $v\in\mathcal B$ and every axis $\alpha$,
\begin{equation}
 \mathcal U_i^{(5)}((s_i(v))_\alpha)
       =\mathcal V_i(v_\alpha),
 \label{eq:ex-trace-spun-response}
\end{equation}
where $\mathcal V_i=\operatorname{Vault}(S,Z,5,k_i,c_i)$ is the grade-five
strip evaluation at scale index $k_i$ with crown normalization $c_i$. For
a raw grade-five label $x$ which is odd
under central signed inversion, meaning
$F(\operatorname{archV})x=-x$ (the night condition), the exact scaling
identity is
\begin{equation}
 \mathcal V_i(\mathcal E_{k_i}(x))
   =c_i P_5x,
 \qquad c_i=\zeta_i(a_i b_i)^5,
 \label{eq:ex-trace-night-scaling}
\end{equation}
where $P_5=S.\operatorname{project}_5$ maps raw grade-five labels into
$S_5$. These are exact
source identities on the stated tail, not consequences of the estimate
below. Every entry $y_\alpha$ of an array $y\in\mathcal B$ satisfies the
night condition: the central inversion fixes each axis and attaches the
sign $-1$ to it, so signed equivariance reads
$F(\operatorname{archV})y_\alpha=-y_\alpha$. In particular
\eqref{eq:ex-trace-night-scaling} applies to $x=x_i(\alpha)$.

The tested configuration consists of two points. With $e_0$ the unit vector
of axis $0$, put
\[
 \zeta_i^{\mathrm{pt}}
   =\bigl(b_i^{-1}(\lceil b_i\rceil,1,1,1),\,0\bigr),
 \qquad
 \zeta_\infty^{\mathrm{pt}}=(e_0,0).
\]
The scale input is $b_i\to\infty$; hence
$\zeta_i^{\mathrm{pt}}\to\zeta_\infty^{\mathrm{pt}}$, the axial
configuration. By $r_i=a_ib_i$, the dilated configuration
$r_i\zeta_i^{\mathrm{pt}}=(a_i(\lceil b_i\rceil,1,1,1),0)$ consists of two
distinct points of the lattice of spacing $a_i$, the kind of configuration
at which the finite-volume price estimate is formulated.
For a reference grade-four field $w\in E_4$ and $Y\in\mathcal V$, let
$\lambda_{i,w}(Y)$ be the
two-point function of the selected table with entries
$\mathcal U_i^{(5)}(Y_0)$ and
$\mathcal U_i^{(4)}w$, at the shrinking configuration
$r_i\zeta_i^{\mathrm{pt}}$; only the axis-$0$ entry $Y_0$ of the array
enters. Let $\lambda_w(Y)$ be the two-point function of the reference table
with entries $Y_0$ and $w$ at $\zeta_\infty^{\mathrm{pt}}$.
Denote the selected-table two-point function at
$r_i\zeta_i^{\mathrm{pt}}$ on explicitly supplied grade-five and grade-four
entries
by $\mathfrak p_{i,w}$; it is complex-linear in the grade-five entry. The
local price input is the bound
\begin{equation}
 \left|\mathfrak p_{i,w}
       (a_i^5P_5x_i(0),\mathcal U_i^{(4)}w)\right|
 \leq \Theta_i,
 \quad\text{eventually in }i\text{ for each fixed }w,
 \label{eq:ex-trace-small-means-price}
\end{equation}
where
\[
 \Theta_i=\exp\!\left(-\frac{100}{10^4}
      (\log\widehat\beta_{j_i})^{48}\right)
\]
is the positive local price allowance
at the frame level $j_i$. The source proves
this receipt from the priced mission and its finite-volume mileage
estimate; Proposition~\ref{prop:ex-trace-mission-price-transfer} is the
corresponding conditional transfer statement, and the bound remains an
upstream input here.

The detector rests on two inputs, which are not proved in this section.
\begin{enumerate}
\item \emph{Separation at the axial configuration.} If $y\in\mathcal V$
and $\lambda_w(y)=0$ for every $w\in E_4$, then $y=0$. In the source this
is derived from the derivative and orientation identities of the reference
table and simplicity of the gauge group, through the injectivity of the
axial divergence probes on the reference vector table; compare
\eqref{eq:ex-trace-divergence-probes} and the isomorphism stated after it.
\item \emph{Convergence of fixed tested pairings.} For fixed $u\in E_5$
and $w\in E_4$, and every sequence of two-point configurations
$\xi_i\to\zeta_\infty^{\mathrm{pt}}$, the selected-table two-point
function with entries $\mathcal U_i^{(5)}u$ and $\mathcal U_i^{(4)}w$ at
$r_i\xi_i$ converges to the reference-table two-point function with
entries $u,w$ at $\zeta_\infty^{\mathrm{pt}}$. This is a property of the
same selected charts $\mathcal U_i^{(d)}$ and dilations $r_i$ as above; in
the source it follows from the convergence of the selected chart path
together with the derivative identities of both tables.
\end{enumerate}
The finite detector is deduced from them as follows. The maps
$y\mapsto\lambda_w(y)$, $w\in E_4$, are real-linear functionals on the
finite-dimensional space $\mathcal V$ with zero common kernel by~(1), so
finitely many of them already have zero common kernel. Hence there are
fixed reference fields
$w_1,\ldots,w_m\in E_4$ such that the limiting pairings define an injective
real-linear map
\[
 \Lambda(y)=(\lambda_{w_r}(y))_{r=1}^m:
       \mathcal V\longrightarrow\mathbb C^m.
\]
By~(2), applied with $\xi_i=\zeta_i^{\mathrm{pt}}$, $u=y_0$ and
$w=w_r$, the scale-dependent real-linear maps
$\Lambda_i(y)=(\lambda_{i,w_r}(y))_{r=1}^m$
converge to $\Lambda$ pointwise on $\mathcal V$, and therefore in operator
norm, since $\mathcal V$ is finite-dimensional. If
$\mathcal V\neq\{0\}$, injectivity gives a
positive minimum $2c$ of $\|\Lambda y\|$ on its unit sphere, and
$\|\Lambda_i-\Lambda\|\le c$ for large $i$. Thus some $c>0$ satisfies
\begin{equation}
 \|\Lambda_i y\|\geq c\|y\|
 \qquad\text{for all sufficiently large }i\text{ and every }y\in\mathcal V.
 \label{eq:ex-trace-small-means-detector}
\end{equation}
For $\mathcal V=\{0\}$ the desired conclusion is immediate. The
separation and convergence inputs, not finite dimensionality alone, are
the detector inputs.

The scalar rate input is
\begin{equation}
 z_i^p m_i^5\Theta_i\longrightarrow0
 \qquad\text{for each fixed integer }p\geq0.
 \label{eq:ex-trace-small-means-rate}
\end{equation}
In the selected source, this follows by multiplying three rate tickets
for sequences indexed by the frame level $j$ and then substituting
$j=j_i\to\infty$. Here a nonnegative sequence $f_j$ has a ticket of rate
$\rho$ if, for every $\rho'>\rho$,
$f_j\le\exp(\rho'(\log\widehat\beta_j)^{48})$ for all sufficiently large
$j$. The three tickets are:
the scale $\widehat\beta$ at the summit index of $j$, which for $j=j_i$
is $z_i$ itself, has a zero-rate ticket (it is eventually at most twice
$\widehat\beta_j$); the cap $\operatorname{tiltcap}(j)$ has a zero-rate
ticket, hence so does its fifth power; and
$\exp(-100\delta(\log\widehat\beta_j)^{48})$, which is $\Theta_i$ for
$j=j_i$,
has rate $-100\delta$, with the fixed source constant $\delta=10^{-4}$.
Ticket products add rates, so for fixed $p$ the product has negative rate
$-100\delta<0$ and tends to zero, because
$\widehat\beta_j\to\infty$. These scalar receipts are listed with
their exact source
declarations in the evidence file.

\begin{lemma}[Weighted field decay from the priced mission]
\label{lem:ex-trace-small-means}
Assume that, for all sufficiently large $i$, the scale relations
$r_i=a_ib_i$ and $0<b_i\le m_i$, the medium representation
\eqref{eq:ex-trace-medium-core}, the response identity
\eqref{eq:ex-trace-spun-response} and the scaling identity
\eqref{eq:ex-trace-night-scaling} for $x=x_i(0)$ hold. Assume the price bound
\eqref{eq:ex-trace-small-means-price} for each of the finitely many
detector fields $w_1,\ldots,w_m$, the detector bound
\eqref{eq:ex-trace-small-means-detector}, and the rate
\eqref{eq:ex-trace-small-means-rate}. Then, for each fixed integer
$p\geq0$,
\[
 z_i^p s_i(v_i)\longrightarrow0\quad\text{in }\mathcal V.
\]
\end{lemma}
\begin{proof}
Set $Y_i=z_i^p s_i(v_i)\in\mathcal V$, and restrict to $i$ so large that
all the assumed identities and the finitely many price bounds hold. By the
medium representation,
response normalization, and night scaling, together with the real
linearity of $s_i$, $\mathcal U_i^{(5)}$ and $\mathcal V_i$,
\[
 \mathcal U_i^{(5)}((Y_i)_0)
 =z_i^p\zeta_i^{-1}c_iP_5x_i(0)
 =z_i^p(a_i b_i)^5P_5x_i(0).
\]
The Arch scale $\zeta_i$ cancels; it is distinct from the polynomial weight
$z_i$. Linearity in the grade-five entry gives, for each detector field
$w_r$,
\[
 \lambda_{i,w_r}(Y_i)
 =z_i^p b_i^5\mathfrak p_{i,w_r}
      (a_i^5P_5x_i(0),\mathcal U_i^{(4)}w_r).
\]
The local price bound and $b_i\leq m_i$ imply
\[
 \|\Lambda_iY_i\|
 \leq \sqrt m\,z_i^p m_i^5\Theta_i\longrightarrow0.
\]
The detector lower bound gives $\|Y_i\|\leq c^{-1}\|\Lambda_iY_i\|\to0$.
This proves the desired limit in the full finite-dimensional response
space.
\end{proof}

The source applies this lemma to its selected medium. The proof identifies
the mechanism: the powered local price controls finitely many separating
coordinates, and the negative rate forces each coordinate to zero. The
uniform local price receipt, mission construction, and detector convergence
still depend on the selected finite-volume funding and chart construction.
